\documentclass[12pt]{article}
\usepackage[utf8]{inputenc}
\usepackage[T1]{fontenc}
\usepackage{hyperref}
\title{Santa In\'es Yatzechi Zapotec Language Resources for K-12 Curriculum Development}
\author{Compiled by Claude Sonnet 5.5 (claude-sonnet-5-5) for Andr\'as Kornai}
\date{}
\begin{document}
\maketitle

\section{Introduction}
Santa In\'es Yatzechi Zapotec (ISO 639-3 zpn, Glottocode sant1447; also called Yatzeche or Zegache Zapotec) is an Otomanguean language of Oaxaca, Mexico, with roughly 2,200 speakers by an older census figure. Search found a short literacy booklet, reference pages, and religious audio. I did not search specifically for dictionaries or grammars, so gaps may remain.

\section{Primary Educational Resources}
\begin{itemize}
\item \emph{C\'omo leer y escribir el zapoteco}, by Donald and Anne Olson (SIL, 1986), a 23-page bilingual Spanish and Zapotec literacy primer for the Santa In\'es Yatzeche variety, held by El Colegio de M\'exico. \url{https://repositorio.colmex.mx/concern/books/9306t3908?locale=en}
\end{itemize}

\section{Linguistic and Reference Materials}
\begin{itemize}
\item Wikipedia article on Yatzeche Zapotec. \url{https://en.wikipedia.org/wiki/Yatzeche_Zapotec}
\item Glottolog entry (sant1447). \url{https://glottolog.org/resource/languoid/id/sant1447.bigmap.html}
\item ASJP word list page, citing a basic vocabulary paper from SIL Mexico. \url{https://asjp.clld.org/languages/SANTA_INES_YATZECHI_ZAPOTEC}
\item Wikitongues language page. \url{https://wikitongues.org/languages/zpn/}
\end{itemize}

\section{Religious and Community Resources}
\begin{itemize}
\item Global Recordings Network: audio and slideshow Bible-teaching and evangelistic materials. \url{https://www.globalrecordings.net/en/language/zpn}
\end{itemize}

\section{Digital and Audio Resources}
The Global Recordings Network materials above are the only audio found, and they are religious in character.

\section{Curriculum Development Context}
The language is not known to be taught in schools. Adults all use it as a first language but not all young people do, so it is endangered. The Olson primer is the only educational item found and is very short. A curriculum would need substantial new materials developed with community speakers.

\end{document}